% Options for packages loaded elsewhere
\PassOptionsToPackage{unicode}{hyperref}
\PassOptionsToPackage{hyphens}{url}
\documentclass[
  11pt,
  a4paper,
]{article}
\usepackage{xcolor}
\usepackage[margin=18mm]{geometry}
\usepackage{amsmath,amssymb}
\setcounter{secnumdepth}{-\maxdimen} % remove section numbering
\usepackage{iftex}
\ifPDFTeX
  \usepackage[T1]{fontenc}
  \usepackage[utf8]{inputenc}
  \usepackage{textcomp} % provide euro and other symbols
\else % if luatex or xetex
  \usepackage{unicode-math} % this also loads fontspec
  \defaultfontfeatures{Scale=MatchLowercase}
  \defaultfontfeatures[\rmfamily]{Ligatures=TeX,Scale=1}
\fi
\ifPDFTeX\else
  % xetex/luatex font selection
    \setmainfont[]{Noto Serif}
    \setsansfont[]{Noto Sans}
    \setmonofont[]{Noto Sans Mono}
  \setmathfont[]{Noto Sans Math}
\fi
% Use upquote if available, for straight quotes in verbatim environments
\IfFileExists{upquote.sty}{\usepackage{upquote}}{}
\IfFileExists{microtype.sty}{% use microtype if available
  \usepackage[]{microtype}
  \UseMicrotypeSet[protrusion]{basicmath} % disable protrusion for tt fonts
}{}
\makeatletter
\@ifundefined{KOMAClassName}{% if non-KOMA class
  \IfFileExists{parskip.sty}{%
    \usepackage{parskip}
  }{% else
    \setlength{\parindent}{0pt}
    \setlength{\parskip}{6pt plus 2pt minus 1pt}}
}{% if KOMA class
  \KOMAoptions{parskip=half}}
\makeatother
\ifLuaTeX
\usepackage[bidi=basic]{babel}
\else
\usepackage[bidi=default]{babel}
\fi
\babelprovide[main,import]{russian}
\ifPDFTeX
\else
\babelfont{rm}[]{Noto Serif}
\fi
% get rid of language-specific shorthands (see #6817):
\let\LanguageShortHands\languageshorthands
\def\languageshorthands#1{}
\setlength{\emergencystretch}{3em} % prevent overfull lines
\providecommand{\tightlist}{%
  \setlength{\itemsep}{0pt}\setlength{\parskip}{0pt}}
\usepackage{bookmark}
\IfFileExists{xurl.sty}{\usepackage{xurl}}{} % add URL line breaks if available
\urlstyle{same}
\hypersetup{
  pdftitle={Билеты по дискретным случайным процессам},
  pdflang={ru-RU},
  hidelinks,
  pdfcreator={LaTeX via pandoc}}

\title{Билеты по дискретным случайным процессам}
\author{}
\date{}

\begin{document}
\maketitle

\section{Билет 11. Замена переменной в интеграле Лебега. Моменты через
интегралы по пространству и по
распределению}\label{ux431ux438ux43bux435ux442-11.-ux437ux430ux43cux435ux43dux430-ux43fux435ux440ux435ux43cux435ux43dux43dux43eux439-ux432-ux438ux43dux442ux435ux433ux440ux430ux43bux435-ux43bux435ux431ux435ux433ux430.-ux43cux43eux43cux435ux43dux442ux44b-ux447ux435ux440ux435ux437-ux438ux43dux442ux435ux433ux440ux430ux43bux44b-ux43fux43e-ux43fux440ux43eux441ux442ux440ux430ux43dux441ux442ux432ux443-ux438-ux43fux43e-ux440ux430ux441ux43fux440ux435ux434ux435ux43bux435ux43dux438ux44e}

Источники: Ширяев 1, гл. II, §6; лекции ДСП.

\subsection{1. Короткий
ответ}\label{ux43aux43eux440ux43eux442ux43aux438ux439-ux43eux442ux432ux435ux442}

Главная идея билета: интеграл функции от случайного элемента можно
считать не по исходному вероятностному пространству, а по распределению
этого элемента.

Если

\[
X:(\Omega,\mathcal F,P)\to(E,\mathcal E)
\]

\begin{itemize}
\tightlist
\item
  случайный элемент, а \(P_X\) - его распределение:
\end{itemize}

\[
P_X(B)=P\{X\in B\},
\qquad B\in\mathcal E,
\]

то для измеримой функции \(g:E\to\mathbb R\) или \(\mathbb C\):

\[
\mathbb E g(X)
=
\int_\Omega g(X(\omega))\,P(d\omega)
=
\int_E g(x)\,P_X(dx),
\]

если \(g\ge0\) или если \(g(X)\) интегрируема.

Если \(X\) имеет плотность \(p_X\) относительно меры Лебега на
\(\mathbb R^n\), то

\[
\mathbb E g(X)=\int_{\mathbb R^n} g(x)p_X(x)\,dx.
\]

Отсюда моменты считаются через распределение:

\[
\mathbb EX^k=\int_{\mathbb R} x^k\,P_X(dx),
\qquad
\mathbb E|X|^k=\int_{\mathbb R} |x|^k\,P_X(dx),
\]

а при плотности:

\[
\mathbb EX^k=\int_{\mathbb R} x^k p_X(x)\,dx.
\]

Для случайного вектора \(X\in\mathbb R^n\):

\[
\mathbb EX=\int_{\mathbb R^n}x\,P_X(dx),
\qquad
\operatorname{Cov}(X)=\int_{\mathbb R^n}(x-m)(x-m)^T\,P_X(dx),
\]

где \(m=\mathbb EX\).

\subsection{2. Полный
ответ}\label{ux43fux43eux43bux43dux44bux439-ux43eux442ux432ux435ux442}

В предыдущих билетах были введены распределения, плотности и интеграл
Лебега. Теперь нужно связать математическое ожидание на вероятностном
пространстве с интегралом по распределению. Это и есть вероятностная
форма замены переменной.

\subsubsection{2.1. Вероятностная форма замены
переменной}\label{ux432ux435ux440ux43eux44fux442ux43dux43eux441ux442ux43dux430ux44f-ux444ux43eux440ux43cux430-ux437ux430ux43cux435ux43dux44b-ux43fux435ux440ux435ux43cux435ux43dux43dux43eux439}

Пусть

\[
(\Omega,\mathcal F,P)
\]

\begin{itemize}
\tightlist
\item
  вероятностное пространство, \((E,\mathcal E)\) - измеримое
  пространство, а
\end{itemize}

\[
X:\Omega\to E
\]

\begin{itemize}
\tightlist
\item
  случайный элемент, то есть \(\mathcal F/\mathcal E\)-измеримое
  отображение. Его распределение, или образ меры \(P\) при отображении
  \(X\), задается формулой
\end{itemize}

\[
P_X(B)=P(X^{-1}(B))=P\{\omega:X(\omega)\in B\},
\qquad B\in\mathcal E.
\]

Распределение \(P_X\) - это вероятностная мера на пространстве значений
\(E\). Оно содержит всю информацию, необходимую для вычисления
вероятностей событий вида \(\{X\in B\}\) и средних величин вида
\(g(X)\).

Теорема о замене переменной в интеграле Лебега для образа меры
формулируется так. Пусть \(g:E\to\mathbb R\) или \(\mathbb C\) -
\(\mathcal E\)-измеримая функция. Тогда

\[
\int_\Omega g(X(\omega))\,P(d\omega)
=
\int_E g(x)\,P_X(dx),
\]

если \(g\ge0\); в этом случае обе части понимаются как расширенные
неотрицательные интегралы. Если \(g\) знакопеременная или
комплекснозначная, то равенство верно при интегрируемости:

\[
\int_E |g(x)|\,P_X(dx)<\infty,
\]

что равносильно

\[
\mathbb E|g(X)|<\infty.
\]

В вероятностной записи:

\[
\mathbb E g(X)=\int_E g(x)\,P_X(dx).
\]

Это не ``формальная замена буквы \(\omega\) на \(x\)''. Смысл в том, что
распределение \(P_X\) уже учитывает, с какими вероятностями значения
\(X\) попадают в разные множества.

Для вещественной случайной величины \(X\) можно писать

\[
\mathbb E g(X)=\int_{\mathbb R} g(x)\,P_X(dx).
\]

Если \(F_X\) - функция распределения, то \(P_X\) является мерой
Лебега-Стилтьеса, порожденной \(F_X\), и часто пишут

\[
\mathbb E g(X)=\int_{\mathbb R} g(x)\,dF_X(x).
\]

Если \(P_X\) абсолютно непрерывно относительно меры Лебега и имеет
плотность \(p_X\), то по Билету 10

\[
P_X(dx)=p_X(x)\,dx,
\]

и формула принимает вид

\[
\mathbb E g(X)=\int_{\mathbb R} g(x)p_X(x)\,dx.
\]

\subsubsection{2.2. Случайные векторы и
процессы}\label{ux441ux43bux443ux447ux430ux439ux43dux44bux435-ux432ux435ux43aux442ux43eux440ux44b-ux438-ux43fux440ux43eux446ux435ux441ux441ux44b}

Для случайного вектора \(X=(X_1,\ldots,X_n)\):

\[
\mathbb E g(X)
=
\int_{\mathbb R^n}g(x_1,\ldots,x_n)\,P_X(dx_1,\ldots,dx_n).
\]

Если есть плотность \(p_X\) на \(\mathbb R^n\), то

\[
\mathbb E g(X)
=
\int_{\mathbb R^n}g(x)p_X(x)\,dx.
\]

Это особенно важно для конечномерных распределений случайного процесса.
Пусть \((X_t)_{t\in T}\) - процесс и выбраны моменты времени

\[
t_1,\ldots,t_n.
\]

Тогда вектор

\[
X^{(t)}=(X_{t_1},\ldots,X_{t_n})
\]

имеет конечномерное распределение \(P_{t_1,\ldots,t_n}\). Для любой
измеримой функции \(g:\mathbb R^n\to\mathbb R\):

\[
\mathbb E g(X_{t_1},\ldots,X_{t_n})
=
\int_{\mathbb R^n}g(x_1,\ldots,x_n)\,P_{t_1,\ldots,t_n}(dx).
\]

Если есть плотность конечномерного распределения \(p_{t_1,\ldots,t_n}\),
то

\[
\mathbb E g(X_{t_1},\ldots,X_{t_n})
=
\int_{\mathbb R^n}g(x_1,\ldots,x_n)
p_{t_1,\ldots,t_n}(x_1,\ldots,x_n)\,dx_1\cdots dx_n.
\]

\subsubsection{2.3. Замена переменной и
диффеоморфизмы}\label{ux437ux430ux43cux435ux43dux430-ux43fux435ux440ux435ux43cux435ux43dux43dux43eux439-ux438-ux434ux438ux444ux444ux435ux43eux43cux43eux440ux444ux438ux437ux43cux44b}

Теперь классическая замена переменной в интеграле Лебега. Пусть
\(U,V\subset\mathbb R^n\) - открытые множества, а

\[
\Phi:U\to V
\]

\begin{itemize}
\tightlist
\item
  взаимно однозначное непрерывно дифференцируемое отображение с
  непрерывно дифференцируемым обратным, то есть \(C^1\)-диффеоморфизм.
  Тогда для неотрицательной или интегрируемой функции \(h\):
\end{itemize}

\[
\int_V h(y)\,dy
=
\int_U h(\Phi(x))\,|\det D\Phi(x)|\,dx.
\]

Эквивалентно, если \(y=\Phi(x)\), то элемент объема меняется по правилу

\[
dy=|\det D\Phi(x)|\,dx.
\]

Эта формула нужна, например, при переходе от декартовых координат к
полярным, при линейных преобразованиях векторов и при нахождении
плотности преобразованной случайной величины.

Если \(Y=\Phi(X)\), где \(X\) имеет плотность \(p_X\) на \(U\), а
\(\Phi:U\to V\) - диффеоморфизм, то плотность \(Y\) равна

\[
p_Y(y)
=
p_X(\Phi^{-1}(y))
\left|\det D\Phi^{-1}(y)\right|.
\]

В одномерном случае, если \(Y=\phi(X)\) и \(\phi\) строго монотонна,
дифференцируема и \(\phi'(x)\ne0\), то

\[
p_Y(y)=p_X(\phi^{-1}(y))
\left|(\phi^{-1})'(y)\right|.
\]

Если отображение не взаимно однозначно, надо суммировать по всем
прообразам:

\[
p_Y(y)=\sum_{x:\phi(x)=y}\frac{p_X(x)}{|\phi'(x)|},
\]

когда такая формула корректна.

\subsubsection{2.4. Моменты и числовые
характеристики}\label{ux43cux43eux43cux435ux43dux442ux44b-ux438-ux447ux438ux441ux43bux43eux432ux44bux435-ux445ux430ux440ux430ux43aux442ux435ux440ux438ux441ux442ux438ux43aux438}

Теперь моменты. Момент порядка \(k\) вещественной случайной величины
\(X\) - это математическое ожидание

\[
\mathbb EX^k,
\]

если оно существует. Абсолютный момент порядка \(k\):

\[
\mathbb E|X|^k.
\]

Существование абсолютного момента означает

\[
\mathbb E|X|^k<\infty.
\]

Через распределение:

\[
\mathbb E|X|^k
=
\int_{\mathbb R} |x|^k\,P_X(dx),
\]

а если этот интеграл конечен, то момент

\[
\mathbb EX^k
=
\int_{\mathbb R} x^k\,P_X(dx)
\]

корректно определен. При плотности:

\[
\mathbb E|X|^k=\int_{\mathbb R} |x|^k p_X(x)\,dx,
\qquad
\mathbb EX^k=\int_{\mathbb R} x^k p_X(x)\,dx.
\]

Математическое ожидание:

\[
\mathbb EX=\int_\Omega X(\omega)\,P(d\omega)
=
\int_{\mathbb R} x\,P_X(dx).
\]

При плотности:

\[
\mathbb EX=\int_{\mathbb R} x p_X(x)\,dx.
\]

Для дискретной случайной величины, принимающей значения \(x_i\) с
вероятностями \(p_i\):

\[
\mathbb E g(X)=\sum_i g(x_i)p_i.
\]

Значит

\[
\mathbb EX=\sum_i x_i p_i,
\qquad
\mathbb EX^k=\sum_i x_i^k p_i.
\]

Дисперсия:

\[
\operatorname{Var}X
=
\mathbb E[(X-\mathbb EX)^2].
\]

Через распределение:

\[
\operatorname{Var}X
=
\int_{\mathbb R} (x-m)^2\,P_X(dx),
\qquad
m=\mathbb EX.
\]

Эквивалентная вычислительная формула:

\[
\operatorname{Var}X=\mathbb EX^2-(\mathbb EX)^2,
\]

если \(\mathbb EX^2<\infty\). При плотности:

\[
\operatorname{Var}X
=
\int_{\mathbb R} (x-m)^2p_X(x)\,dx
=
\int_{\mathbb R} x^2p_X(x)\,dx
-
\left(\int_{\mathbb R} xp_X(x)\,dx\right)^2.
\]

\subsubsection{2.5. Ковариация и характеристики случайных
векторов}\label{ux43aux43eux432ux430ux440ux438ux430ux446ux438ux44f-ux438-ux445ux430ux440ux430ux43aux442ux435ux440ux438ux441ux442ux438ux43aux438-ux441ux43bux443ux447ux430ux439ux43dux44bux445-ux432ux435ux43aux442ux43eux440ux43eux432}

Для двух случайных величин \(X,Y\) с конечными вторыми моментами
ковариация:

\[
\operatorname{Cov}(X,Y)
=
\mathbb E[(X-\mathbb EX)(Y-\mathbb EY)].
\]

Если \((X,Y)\) имеет совместное распределение \(P_{X,Y}\), то

\[
\operatorname{Cov}(X,Y)
=
\int_{\mathbb R^2}(x-m_X)(y-m_Y)\,P_{X,Y}(dx,dy).
\]

При совместной плотности \(p_{X,Y}\):

\[
\operatorname{Cov}(X,Y)
=
\int_{\mathbb R^2}(x-m_X)(y-m_Y)p_{X,Y}(x,y)\,dx\,dy.
\]

Также

\[
\operatorname{Cov}(X,Y)=\mathbb E[XY]-\mathbb EX\,\mathbb EY,
\]

если \(\mathbb E X^2<\infty\) и \(\mathbb E Y^2<\infty\).

Для случайного вектора \(X=(X_1,\ldots,X_n)^T\) с конечным вторым
моментом:

\[
m=\mathbb EX
=
\int_{\mathbb R^n}x\,P_X(dx),
\]

а матрица вторых моментов:

\[
M=\mathbb E[XX^T]
=
\int_{\mathbb R^n}xx^T\,P_X(dx).
\]

Ковариационная матрица:

\[
\Sigma
=
\operatorname{Cov}(X)
=
\mathbb E[(X-m)(X-m)^T]
=
\int_{\mathbb R^n}(x-m)(x-m)^T\,P_X(dx).
\]

Эквивалентно:

\[
\Sigma=M-mm^T.
\]

При плотности \(p_X\) все интегралы по \(P_X\) заменяются на интегралы
по мере Лебега:

\[
m=\int_{\mathbb R^n}x p_X(x)\,dx,
\]

\[
M=\int_{\mathbb R^n}xx^T p_X(x)\,dx,
\]

\[
\Sigma=\int_{\mathbb R^n}(x-m)(x-m)^T p_X(x)\,dx.
\]

Элементы матрицы:

\[
\Sigma_{ij}
=
\operatorname{Cov}(X_i,X_j)
=
\int_{\mathbb R^n}(x_i-m_i)(x_j-m_j)\,P_X(dx).
\]

Если есть плотность:

\[
\Sigma_{ij}
=
\int_{\mathbb R^n}(x_i-m_i)(x_j-m_j)p_X(x)\,dx.
\]

Ковариационная матрица всегда симметрична и неотрицательно определена:

\[
a^T\Sigma a
=
\operatorname{Var}(a^TX)
\ge0
\]

для любого \(a\in\mathbb R^n\). Это связывает билет с Билетом 14 и
Билетом 18.

Итак, основная логика билета: случайная величина живет на \(\Omega\), но
ее распределение живет на пространстве значений. Теорема о замене
переменной позволяет считать все средние функции от случайной величины
через распределение. Если есть плотность, интеграл по распределению
превращается в обычный лебегов интеграл с плотностью.

\subsection{3. Основные
определения}\label{ux43eux441ux43dux43eux432ux43dux44bux435-ux43eux43fux440ux435ux434ux435ux43bux435ux43dux438ux44f}

\subsubsection{Образ
меры}\label{ux43eux431ux440ux430ux437-ux43cux435ux440ux44b}

Если \(X:(\Omega,\mathcal F)\to(E,\mathcal E)\) измеримо, а \(P\) - мера
на \(\Omega\), то образ меры \(P\) при \(X\) задается формулой

\[
P_X(B)=P(X^{-1}(B)),
\qquad B\in\mathcal E.
\]

В вероятности \(P_X\) называется распределением случайного элемента
\(X\).

\subsubsection{Замена переменной через образ
меры}\label{ux437ux430ux43cux435ux43dux430-ux43fux435ux440ux435ux43cux435ux43dux43dux43eux439-ux447ux435ux440ux435ux437-ux43eux431ux440ux430ux437-ux43cux435ux440ux44b}

Для измеримой \(g\):

\[
\int_\Omega g(X(\omega))\,P(d\omega)
=
\int_E g(x)\,P_X(dx).
\]

\subsubsection{Интеграл Лебега-Стилтьеса по функции
распределения}\label{ux438ux43dux442ux435ux433ux440ux430ux43b-ux43bux435ux431ux435ux433ux430-ux441ux442ux438ux43bux442ux44cux435ux441ux430-ux43fux43e-ux444ux443ux43dux43aux446ux438ux438-ux440ux430ux441ux43fux440ux435ux434ux435ux43bux435ux43dux438ux44f}

Если \(X\) вещественная, \(F_X\) - ее функция распределения, а \(P_X\) -
соответствующая мера, то

\[
\int_{\mathbb R} g(x)\,P_X(dx)
\]

часто записывают как

\[
\int_{\mathbb R} g(x)\,dF_X(x).
\]

\subsubsection{Плотность
распределения}\label{ux43fux43bux43eux442ux43dux43eux441ux442ux44c-ux440ux430ux441ux43fux440ux435ux434ux435ux43bux435ux43dux438ux44f}

Если \(P_X\ll\lambda_n\), то

\[
p_X=\frac{dP_X}{d\lambda_n}
\]

и

\[
P_X(B)=\int_B p_X(x)\,dx.
\]

\subsubsection{Классическая замена переменной с
якобианом}\label{ux43aux43bux430ux441ux441ux438ux447ux435ux441ux43aux430ux44f-ux437ux430ux43cux435ux43dux430-ux43fux435ux440ux435ux43cux435ux43dux43dux43eux439-ux441-ux44fux43aux43eux431ux438ux430ux43dux43eux43c}

Если \(\Phi:U\to V\) - \(C^1\)-диффеоморфизм, то

\[
\int_V h(y)\,dy
=
\int_U h(\Phi(x))|\det D\Phi(x)|\,dx.
\]

\subsubsection{\texorpdfstring{Момент порядка
\(k\)}{Момент порядка k}}\label{ux43cux43eux43cux435ux43dux442-ux43fux43eux440ux44fux434ux43aux430-k}

Для вещественной случайной величины:

\[
\mathbb EX^k,
\]

если математическое ожидание существует.

\subsubsection{\texorpdfstring{Абсолютный момент порядка
\(k\)}{Абсолютный момент порядка k}}\label{ux430ux431ux441ux43eux43bux44eux442ux43dux44bux439-ux43cux43eux43cux435ux43dux442-ux43fux43eux440ux44fux434ux43aux430-k}

\[
\mathbb E|X|^k.
\]

Если \(\mathbb E|X|^k<\infty\), то момент \(\mathbb EX^k\) существует.

\subsubsection{Дисперсия}\label{ux434ux438ux441ux43fux435ux440ux441ux438ux44f}

\[
\operatorname{Var}X=\mathbb E[(X-\mathbb EX)^2],
\]

если \(\mathbb EX^2<\infty\).

\subsubsection{Ковариация}\label{ux43aux43eux432ux430ux440ux438ux430ux446ux438ux44f}

\[
\operatorname{Cov}(X,Y)
=
\mathbb E[(X-\mathbb EX)(Y-\mathbb EY)],
\]

если \(X,Y\in L^2\).

\subsubsection{Матрица вторых
моментов}\label{ux43cux430ux442ux440ux438ux446ux430-ux432ux442ux43eux440ux44bux445-ux43cux43eux43cux435ux43dux442ux43eux432}

\[
M=\mathbb E[XX^T].
\]

\subsubsection{Ковариационная
матрица}\label{ux43aux43eux432ux430ux440ux438ux430ux446ux438ux43eux43dux43dux430ux44f-ux43cux430ux442ux440ux438ux446ux430}

\[
\Sigma=\mathbb E[(X-\mathbb EX)(X-\mathbb EX)^T].
\]

\subsection{4. Основные теоремы и
свойства}\label{ux43eux441ux43dux43eux432ux43dux44bux435-ux442ux435ux43eux440ux435ux43cux44b-ux438-ux441ux432ux43eux439ux441ux442ux432ux430}

\subsubsection{Теорема 1. Замена переменной для образа
меры}\label{ux442ux435ux43eux440ux435ux43cux430-1.-ux437ux430ux43cux435ux43dux430-ux43fux435ux440ux435ux43cux435ux43dux43dux43eux439-ux434ux43bux44f-ux43eux431ux440ux430ux437ux430-ux43cux435ux440ux44b}

Пусть \(X:(\Omega,\mathcal F,P)\to(E,\mathcal E)\) - случайный элемент,
\(P_X\) - его распределение, \(g\) - измеримая функция. Тогда

\[
\int_\Omega g(X(\omega))\,P(d\omega)
=
\int_E g(x)\,P_X(dx),
\]

если \(g\ge0\) или если один из интегралов от \(|g|\) конечен.

\subsubsection{Теорема 2. Доказательная схема замены
переменной}\label{ux442ux435ux43eux440ux435ux43cux430-2.-ux434ux43eux43aux430ux437ux430ux442ux435ux43bux44cux43dux430ux44f-ux441ux445ux435ux43cux430-ux437ux430ux43cux435ux43dux44b-ux43fux435ux440ux435ux43cux435ux43dux43dux43eux439}

Сначала формула проверяется для индикатора:

\[
g=\mathbf 1_B.
\]

Тогда

\[
\int_\Omega \mathbf 1_B(X(\omega))\,dP
=
P\{X\in B\}
=
P_X(B)
=
\int_E\mathbf 1_B(x)\,P_X(dx).
\]

Затем формула переносится на простые функции по линейности, на
неотрицательные измеримые функции по теореме Беппо Леви, а на
интегрируемые знакопеременные функции через положительную и
отрицательную части.

\subsubsection{Теорема 3. Случай с
плотностью}\label{ux442ux435ux43eux440ux435ux43cux430-3.-ux441ux43bux443ux447ux430ux439-ux441-ux43fux43bux43eux442ux43dux43eux441ux442ux44cux44e}

Если \(X\) имеет плотность \(p_X\) относительно \(\lambda_n\), то

\[
\mathbb E g(X)=\int_{\mathbb R^n}g(x)p_X(x)\,dx,
\]

для \(g\ge0\) или интегрируемой \(g(X)\).

\subsubsection{\texorpdfstring{Теорема 4. Классическая замена переменной
в
\(\mathbb R^n\)}{Теорема 4. Классическая замена переменной в \textbackslash mathbb R\^{}n}}\label{ux442ux435ux43eux440ux435ux43cux430-4.-ux43aux43bux430ux441ux441ux438ux447ux435ux441ux43aux430ux44f-ux437ux430ux43cux435ux43dux430-ux43fux435ux440ux435ux43cux435ux43dux43dux43eux439-ux432-mathbb-rn}

Если \(\Phi:U\to V\) - \(C^1\)-диффеоморфизм, то

\[
\int_V h(y)\,dy
=
\int_U h(\Phi(x))|\det D\Phi(x)|\,dx.
\]

\subsubsection{Следствие 5. Плотность
преобразования}\label{ux441ux43bux435ux434ux441ux442ux432ux438ux435-5.-ux43fux43bux43eux442ux43dux43eux441ux442ux44c-ux43fux440ux435ux43eux431ux440ux430ux437ux43eux432ux430ux43dux438ux44f}

Если \(Y=\Phi(X)\), \(X\) имеет плотность \(p_X\), а \(\Phi\) -
диффеоморфизм, то

\[
p_Y(y)=p_X(\Phi^{-1}(y))
\left|\det D\Phi^{-1}(y)\right|.
\]

\subsubsection{Свойство 6. Моменты через
распределение}\label{ux441ux432ux43eux439ux441ux442ux432ux43e-6.-ux43cux43eux43cux435ux43dux442ux44b-ux447ux435ux440ux435ux437-ux440ux430ux441ux43fux440ux435ux434ux435ux43bux435ux43dux438ux435}

Для вещественной \(X\):

\[
\mathbb E|X|^k=\int_{\mathbb R} |x|^k\,P_X(dx),
\]

\[
\mathbb EX^k=\int_{\mathbb R} x^k\,P_X(dx),
\]

если соответствующие интегралы существуют.

\subsubsection{Свойство 7. Моменты при
плотности}\label{ux441ux432ux43eux439ux441ux442ux432ux43e-7.-ux43cux43eux43cux435ux43dux442ux44b-ux43fux440ux438-ux43fux43bux43eux442ux43dux43eux441ux442ux438}

Если \(X\) имеет плотность \(p_X\), то

\[
\mathbb E|X|^k=\int_{\mathbb R} |x|^k p_X(x)\,dx,
\]

\[
\mathbb EX^k=\int_{\mathbb R} x^k p_X(x)\,dx.
\]

\subsubsection{Свойство 8. Дисперсия через
распределение}\label{ux441ux432ux43eux439ux441ux442ux432ux43e-8.-ux434ux438ux441ux43fux435ux440ux441ux438ux44f-ux447ux435ux440ux435ux437-ux440ux430ux441ux43fux440ux435ux434ux435ux43bux435ux43dux438ux435}

Если \(m=\mathbb EX\) и \(\mathbb EX^2<\infty\), то

\[
\operatorname{Var}X=\int_{\mathbb R} (x-m)^2\,P_X(dx)
=
\mathbb EX^2-m^2.
\]

При плотности:

\[
\operatorname{Var}X=\int_{\mathbb R} (x-m)^2p_X(x)\,dx.
\]

\subsubsection{Свойство 9. Ковариация через совместное
распределение}\label{ux441ux432ux43eux439ux441ux442ux432ux43e-9.-ux43aux43eux432ux430ux440ux438ux430ux446ux438ux44f-ux447ux435ux440ux435ux437-ux441ux43eux432ux43cux435ux441ux442ux43dux43eux435-ux440ux430ux441ux43fux440ux435ux434ux435ux43bux435ux43dux438ux435}

Если \(X,Y\in L^2\), то

\[
\operatorname{Cov}(X,Y)
=
\int_{\mathbb R^2}(x-m_X)(y-m_Y)\,P_{X,Y}(dx,dy).
\]

При совместной плотности:

\[
\operatorname{Cov}(X,Y)
=
\int_{\mathbb R^2}(x-m_X)(y-m_Y)p_{X,Y}(x,y)\,dx\,dy.
\]

\subsubsection{Свойство 10. Ковариационная
матрица}\label{ux441ux432ux43eux439ux441ux442ux432ux43e-10.-ux43aux43eux432ux430ux440ux438ux430ux446ux438ux43eux43dux43dux430ux44f-ux43cux430ux442ux440ux438ux446ux430}

Для \(X\in\mathbb R^n\) с конечным вторым моментом:

\[
\Sigma=\int_{\mathbb R^n}(x-m)(x-m)^T\,P_X(dx).
\]

При плотности:

\[
\Sigma=\int_{\mathbb R^n}(x-m)(x-m)^T p_X(x)\,dx.
\]

\subsubsection{Свойство 11. Неотрицательная определенность
ковариационной
матрицы}\label{ux441ux432ux43eux439ux441ux442ux432ux43e-11.-ux43dux435ux43eux442ux440ux438ux446ux430ux442ux435ux43bux44cux43dux430ux44f-ux43eux43fux440ux435ux434ux435ux43bux435ux43dux43dux43eux441ux442ux44c-ux43aux43eux432ux430ux440ux438ux430ux446ux438ux43eux43dux43dux43eux439-ux43cux430ux442ux440ux438ux446ux44b}

Для любого \(a\in\mathbb R^n\):

\[
a^T\Sigma a=\operatorname{Var}(a^TX)\ge0.
\]

\subsection{5. Примеры}\label{ux43fux440ux438ux43cux435ux440ux44b}

\subsubsection{Пример 1. Дискретная случайная
величина}\label{ux43fux440ux438ux43cux435ux440-1.-ux434ux438ux441ux43aux440ux435ux442ux43dux430ux44f-ux441ux43bux443ux447ux430ux439ux43dux430ux44f-ux432ux435ux43bux438ux447ux438ux43dux430}

Пусть \(X\) принимает значения \(x_i\) с вероятностями \(p_i\). Тогда
распределение

\[
P_X=\sum_i p_i\delta_{x_i}.
\]

Поэтому

\[
\mathbb E g(X)=\int g\,dP_X=\sum_i g(x_i)p_i.
\]

В частности:

\[
\mathbb EX=\sum_i x_i p_i,
\qquad
\operatorname{Var}X=\sum_i(x_i-m)^2p_i.
\]

\subsubsection{Пример 2. Непрерывная случайная величина с
плотностью}\label{ux43fux440ux438ux43cux435ux440-2.-ux43dux435ux43fux440ux435ux440ux44bux432ux43dux430ux44f-ux441ux43bux443ux447ux430ux439ux43dux430ux44f-ux432ux435ux43bux438ux447ux438ux43dux430-ux441-ux43fux43bux43eux442ux43dux43eux441ux442ux44cux44e}

Пусть \(X\) имеет плотность

\[
p_X(x)=\mathbf 1_{[0,1]}(x).
\]

Тогда

\[
\mathbb EX=\int_0^1 x\,dx=\frac12,
\]

\[
\mathbb EX^2=\int_0^1 x^2\,dx=\frac13,
\]

\[
\operatorname{Var}X=\frac13-\frac14=\frac1{12}.
\]

\subsubsection{Пример 3. Вектор с
плотностью}\label{ux43fux440ux438ux43cux435ux440-3.-ux432ux435ux43aux442ux43eux440-ux441-ux43fux43bux43eux442ux43dux43eux441ux442ux44cux44e}

Пусть \((X,Y)\) имеет совместную плотность \(p(x,y)\). Тогда

\[
\mathbb E[XY]
=
\int_{\mathbb R^2}xy\,p(x,y)\,dx\,dy.
\]

Если \(m_X=\mathbb EX\), \(m_Y=\mathbb EY\), то

\[
\operatorname{Cov}(X,Y)
=
\int_{\mathbb R^2}(x-m_X)(y-m_Y)p(x,y)\,dx\,dy.
\]

\subsubsection{Пример 4. Полярные
координаты}\label{ux43fux440ux438ux43cux435ux440-4.-ux43fux43eux43bux44fux440ux43dux44bux435-ux43aux43eux43eux440ux434ux438ux43dux430ux442ux44b}

Для перехода

\[
x=r\cos\theta,\qquad y=r\sin\theta
\]

якобиан равен

\[
\left|\det\frac{\partial(x,y)}{\partial(r,\theta)}\right|=r.
\]

Поэтому

\[
\int_{\mathbb R^2}h(x,y)\,dx\,dy
=
\int_0^\infty\int_0^{2\pi}
h(r\cos\theta,r\sin\theta)\,r\,d\theta\,dr.
\]

\subsubsection{Пример 5. Линейное преобразование
плотности}\label{ux43fux440ux438ux43cux435ux440-5.-ux43bux438ux43dux435ux439ux43dux43eux435-ux43fux440ux435ux43eux431ux440ux430ux437ux43eux432ux430ux43dux438ux435-ux43fux43bux43eux442ux43dux43eux441ux442ux438}

Пусть \(X\in\mathbb R^n\) имеет плотность \(p_X\), а

\[
Y=AX+b,
\]

где \(A\) - невырожденная матрица. Тогда

\[
p_Y(y)
=
p_X(A^{-1}(y-b))|\det A^{-1}|
=
\frac{p_X(A^{-1}(y-b))}{|\det A|}.
\]

Для моментов:

\[
\mathbb EY=A\mathbb EX+b,
\]

\[
\operatorname{Cov}(Y)=A\operatorname{Cov}(X)A^T.
\]

\subsubsection{Пример 6. Момент может не
существовать}\label{ux43fux440ux438ux43cux435ux440-6.-ux43cux43eux43cux435ux43dux442-ux43cux43eux436ux435ux442-ux43dux435-ux441ux443ux449ux435ux441ux442ux432ux43eux432ux430ux442ux44c}

У распределения Коши с плотностью

\[
p(x)=\frac1{\pi(1+x^2)}
\]

математическое ожидание не существует как конечный лебегов интеграл,
потому что

\[
\int_{\mathbb R} |x|p(x)\,dx=\infty.
\]

Поэтому нельзя автоматически писать \(\mathbb EX=\int xp(x)\,dx\) как
конечное число.

\subsection{6. Схемы доказательств /
обоснований}\label{ux441ux445ux435ux43cux44b-ux434ux43eux43aux430ux437ux430ux442ux435ux43bux44cux441ux442ux432-ux43eux431ux43eux441ux43dux43eux432ux430ux43dux438ux439}

\textbf{Доказательство формулы \(\mathbb E g(X)=\int g\,dP_X\).}

\begin{enumerate}
\def\labelenumi{\arabic{enumi}.}
\tightlist
\item
  Для \(g=\mathbf 1_B\):
\end{enumerate}

\[
\mathbb E\mathbf 1_B(X)=P\{X\in B\}=P_X(B).
\]

\begin{enumerate}
\def\labelenumi{\arabic{enumi}.}
\setcounter{enumi}{1}
\tightlist
\item
  Для простой функции
\end{enumerate}

\[
g=\sum_{k=1}^n a_k\mathbf 1_{B_k}
\]

используем линейность интеграла.

\begin{enumerate}
\def\labelenumi{\arabic{enumi}.}
\setcounter{enumi}{2}
\item
  Для \(g\ge0\) берем простые функции \(g_n\uparrow g\) и применяем
  теорему Беппо Леви.
\item
  Для знакопеременной \(g\) пишем
\end{enumerate}

\[
g=g^+-g^-,
\]

и требуем интегрируемость, чтобы не получить неопределенность
\(\infty-\infty\).

\textbf{Обоснование формулы с плотностью.}

Если

\[
P_X(dx)=p_X(x)\,dx,
\]

то

\[
\int g\,dP_X
=
\int g(x)p_X(x)\,dx
\]

по определению плотности как производной Радона-Никодима.

\textbf{Обоснование дисперсии.}

Если \(m=\mathbb EX\), то

\[
\operatorname{Var}X
=
\mathbb E[(X-m)^2]
=
\mathbb E[X^2-2mX+m^2]
\]

\[
=
\mathbb EX^2-2m\mathbb EX+m^2
=
\mathbb EX^2-m^2.
\]

\textbf{Обоснование неотрицательной определенности ковариационной
матрицы.}

Для любого \(a\in\mathbb R^n\):

\[
a^T\Sigma a
=
\mathbb E[a^T(X-m)(X-m)^Ta]
=
\mathbb E[(a^T(X-m))^2]
\ge0.
\]

\textbf{Обоснование плотности преобразования.}

Если \(Y=\Phi(X)\), то для борелевского \(B\subset V\):

\[
P\{Y\in B\}
=
\int_{\Phi^{-1}(B)}p_X(x)\,dx.
\]

По классической замене переменной \(x=\Phi^{-1}(y)\):

\[
\int_{\Phi^{-1}(B)}p_X(x)\,dx
=
\int_B p_X(\Phi^{-1}(y))
\left|\det D\Phi^{-1}(y)\right|\,dy.
\]

Значит это и есть плотность \(Y\).

\subsection{7. Что написать на
доске}\label{ux447ux442ux43e-ux43dux430ux43fux438ux441ux430ux442ux44c-ux43dux430-ux434ux43eux441ux43aux435}

\begin{enumerate}
\def\labelenumi{\arabic{enumi}.}
\tightlist
\item
  Образ меры:
\end{enumerate}

\[
P_X(B)=P\{X\in B\}.
\]

\begin{enumerate}
\def\labelenumi{\arabic{enumi}.}
\setcounter{enumi}{1}
\tightlist
\item
  Замена переменной:
\end{enumerate}

\[
\mathbb E g(X)
=
\int_\Omega g(X(\omega))\,dP
=
\int_E g(x)\,P_X(dx).
\]

\begin{enumerate}
\def\labelenumi{\arabic{enumi}.}
\setcounter{enumi}{2}
\tightlist
\item
  При плотности:
\end{enumerate}

\[
\mathbb E g(X)
=
\int_{\mathbb R^n}g(x)p_X(x)\,dx.
\]

\begin{enumerate}
\def\labelenumi{\arabic{enumi}.}
\setcounter{enumi}{3}
\tightlist
\item
  Классическая замена:
\end{enumerate}

\[
\int_V h(y)\,dy
=
\int_U h(\Phi(x))|\det D\Phi(x)|\,dx.
\]

\begin{enumerate}
\def\labelenumi{\arabic{enumi}.}
\setcounter{enumi}{4}
\tightlist
\item
  Моменты:
\end{enumerate}

\[
\mathbb EX^k=\int x^k\,P_X(dx),
\qquad
\mathbb E|X|^k=\int |x|^k\,P_X(dx).
\]

\begin{enumerate}
\def\labelenumi{\arabic{enumi}.}
\setcounter{enumi}{5}
\tightlist
\item
  Дисперсия:
\end{enumerate}

\[
\operatorname{Var}X
=
\int (x-m)^2\,P_X(dx)
=
\mathbb EX^2-m^2.
\]

\begin{enumerate}
\def\labelenumi{\arabic{enumi}.}
\setcounter{enumi}{6}
\tightlist
\item
  Ковариационная матрица:
\end{enumerate}

\[
\Sigma=\int (x-m)(x-m)^T\,P_X(dx),
\qquad
a^T\Sigma a\ge0.
\]

\subsection{8. Типичные дополнительные
вопросы}\label{ux442ux438ux43fux438ux447ux43dux44bux435-ux434ux43eux43fux43eux43bux43dux438ux442ux435ux43bux44cux43dux44bux435-ux432ux43eux43fux440ux43eux441ux44b}

\textbf{Вопрос. Что такое образ меры?}

Это мера на пространстве значений:

\[
P_X(B)=P(X^{-1}(B)).
\]

В вероятности она называется распределением \(X\).

\textbf{Вопрос. В чем смысл формулы \(\mathbb E g(X)=\int g\,dP_X\)?}

Она говорит, что для вычисления среднего функции от \(X\) не нужно знать
все вероятностное пространство. Достаточно знать распределение \(X\).

\textbf{Вопрос. При каких условиях верна формула замены переменной?}

Для \(g\ge0\) она верна в расширенном смысле. Для знакопеременной
функции нужна интегрируемость:

\[
\int |g|\,dP_X<\infty.
\]

\textbf{Вопрос. Чем интеграл по распределению отличается от интеграла по
плотности?}

Интеграл по распределению

\[
\int g\,dP_X
\]

имеет смысл всегда. Интеграл

\[
\int g(x)p_X(x)\,dx
\]

требует существования плотности \(p_X\) относительно меры Лебега.

\textbf{Вопрос. Что делать, если плотности нет?}

Интегрировать по мере распределения \(P_X\) или использовать сумму в
дискретном случае. Нельзя искусственно писать лебегову плотность.

\textbf{Вопрос. Как посчитать момент порядка \(k\)?}

Через замену переменной с \(g(x)=x^k\):

\[
\mathbb EX^k=\int x^k\,P_X(dx),
\]

если интеграл корректно определен.

\textbf{Вопрос. Чем момент отличается от абсолютного момента?}

Абсолютный момент:

\[
\mathbb E|X|^k.
\]

Если он конечен, то обычный момент \(\mathbb EX^k\) точно существует.
Обратное в смысле условной сходимости для Лебега не используется:
математическое ожидание требует интегрируемости положительной и
отрицательной частей без неопределенности.

\textbf{Вопрос. Как выразить дисперсию через распределение?}

Если \(m=\mathbb EX\), то

\[
\operatorname{Var}X=\int (x-m)^2\,P_X(dx).
\]

При плотности:

\[
\operatorname{Var}X=\int (x-m)^2p_X(x)\,dx.
\]

\textbf{Вопрос. Как выразить ковариацию через совместное распределение?}

\[
\operatorname{Cov}(X,Y)
=
\int (x-m_X)(y-m_Y)\,P_{X,Y}(dx,dy).
\]

При совместной плотности заменить \(P_{X,Y}(dx,dy)\) на
\(p_{X,Y}(x,y)\,dx\,dy\).

\textbf{Вопрос. Почему ковариационная матрица неотрицательно
определена?}

Потому что

\[
a^T\Sigma a=\operatorname{Var}(a^TX)\ge0.
\]

\textbf{Вопрос. Как меняется плотность при диффеоморфизме?}

Если \(Y=\Phi(X)\), то

\[
p_Y(y)=p_X(\Phi^{-1}(y))
\left|\det D\Phi^{-1}(y)\right|.
\]

\textbf{Вопрос. Почему нельзя забывать якобиан?}

Потому что мера Лебега объема меняется при растяжении координат. Якобиан
учитывает локальное изменение объема.

\subsection{9. Типичные
ошибки}\label{ux442ux438ux43fux438ux447ux43dux44bux435-ux43eux448ux438ux431ux43aux438}

\begin{enumerate}
\def\labelenumi{\arabic{enumi}.}
\tightlist
\item
  Писать интеграл по плотности, когда плотности нет.
\item
  Путать \(P\) на \(\Omega\) и \(P_X\) на пространстве значений.
\item
  Забывать условия интегрируемости для моментов.
\item
  Считать, что \(\int xp(x)\,dx\) существует как число без проверки
  \(\int |x|p(x)\,dx<\infty\).
\item
  Путать момент \(\mathbb EX^2\) и дисперсию \(\operatorname{Var}X\).
\item
  Забывать центрирование в ковариации.
\item
  Считать ковариационную матрицу матрицей вторых моментов без вычитания
  \(mm^T\).
\item
  Забывать якобиан в классической замене переменной.
\item
  В формуле плотности преобразования использовать \(|\det D\Phi|\)
  вместо \(|\det D\Phi^{-1}|\) при записи через \(y\).
\item
  Интегрировать по маргинальному распределению вместо совместного при
  вычислении \(\mathbb E g(X,Y)\).
\end{enumerate}

\subsection{10. Связи с другими
билетами}\label{ux441ux432ux44fux437ux438-ux441-ux434ux440ux443ux433ux438ux43cux438-ux431ux438ux43bux435ux442ux430ux43cux438}

\textbf{Билет 06.} Распределения случайных величин, векторов и
конечномерные распределения - это меры, по которым интегрируют в этом
билете.

\textbf{Билет 07.} Все формулы основаны на интеграле Лебега и предельных
теоремах для перехода от простых функций к измеримым.

\textbf{Билет 08.} Характеристическая функция - это частный случай:

\[
\varphi_X(t)=\mathbb E e^{itX}
=
\int e^{itx}\,P_X(dx).
\]

\textbf{Билет 09.} Для совместных распределений и плотностей
используются Фубини и повторные интегралы.

\textbf{Билет 10.} Плотность \(p_X\) - это производная Радона-Никодима
\(dP_X/d\lambda\).

\textbf{Билет 13.} Моменты можно выражать через производные
характеристической функции, а здесь - через интегралы по распределению.

\textbf{Билет 14.} Ковариация и корреляция строятся из вторых моментов.

\textbf{Билет 18.} Гауссовский закон задается средним вектором и
ковариационной матрицей.

\textbf{Билет 24.} Корреляционная функция стационарной
последовательности выражается через моменты второго порядка.

\subsection{11. Минимум на
удовлетворительно}\label{ux43cux438ux43dux438ux43cux443ux43c-ux43dux430-ux443ux434ux43eux432ux43bux435ux442ux432ux43eux440ux438ux442ux435ux43bux44cux43dux43e}

Нужно уметь сказать:

\begin{enumerate}
\def\labelenumi{\arabic{enumi}.}
\tightlist
\item
  Распределение:
\end{enumerate}

\[
P_X(B)=P\{X\in B\}.
\]

\begin{enumerate}
\def\labelenumi{\arabic{enumi}.}
\setcounter{enumi}{1}
\tightlist
\item
  Замена переменной:
\end{enumerate}

\[
\mathbb E g(X)
=
\int_\Omega g(X)\,dP
=
\int_E g(x)\,P_X(dx).
\]

\begin{enumerate}
\def\labelenumi{\arabic{enumi}.}
\setcounter{enumi}{2}
\tightlist
\item
  При плотности:
\end{enumerate}

\[
\mathbb E g(X)=\int g(x)p_X(x)\,dx.
\]

\begin{enumerate}
\def\labelenumi{\arabic{enumi}.}
\setcounter{enumi}{3}
\tightlist
\item
  Момент:
\end{enumerate}

\[
\mathbb EX^k=\int x^k\,P_X(dx).
\]

\begin{enumerate}
\def\labelenumi{\arabic{enumi}.}
\setcounter{enumi}{4}
\tightlist
\item
  Дисперсия:
\end{enumerate}

\[
\operatorname{Var}X=\mathbb EX^2-(\mathbb EX)^2.
\]

\subsection{12. Минимум на
хорошо}\label{ux43cux438ux43dux438ux43cux443ux43c-ux43dux430-ux445ux43eux440ux43eux448ux43e}

Помимо минимума, нужно:

\begin{enumerate}
\def\labelenumi{\arabic{enumi}.}
\tightlist
\item
  объяснить, что формула верна для \(g\ge0\) или интегрируемой \(g\);
\item
  различать интеграл по \(P_X\) и интеграл по плотности;
\item
  записать ковариацию через совместное распределение;
\item
  записать матрицу
\end{enumerate}

\[
\Sigma=\int (x-m)(x-m)^T\,P_X(dx);
\]

\begin{enumerate}
\def\labelenumi{\arabic{enumi}.}
\setcounter{enumi}{4}
\tightlist
\item
  привести дискретный и непрерывный пример;
\item
  назвать классическую замену переменной с якобианом.
\end{enumerate}

\subsection{13. Что добавить для
отлично}\label{ux447ux442ux43e-ux434ux43eux431ux430ux432ux438ux442ux44c-ux434ux43bux44f-ux43eux442ux43bux438ux447ux43dux43e}

Для сильного ответа стоит добавить:

\begin{enumerate}
\def\labelenumi{\arabic{enumi}.}
\tightlist
\item
  доказательство через индикаторы, простые функции и теорему Беппо Леви;
\item
  интеграл Лебега-Стилтьеса
\end{enumerate}

\[
\int g\,dF_X;
\]

\begin{enumerate}
\def\labelenumi{\arabic{enumi}.}
\setcounter{enumi}{2}
\tightlist
\item
  формулу плотности преобразования:
\end{enumerate}

\[
p_Y(y)=p_X(\Phi^{-1}(y))|\det D\Phi^{-1}(y)|;
\]

\begin{enumerate}
\def\labelenumi{\arabic{enumi}.}
\setcounter{enumi}{3}
\tightlist
\item
  неотрицательную определенность ковариационной матрицы;
\item
  предупреждение про распределение Коши или другой пример, где момент не
  существует;
\item
  связь с характеристическими функциями и гауссовскими законами.
\end{enumerate}

\subsection{14. Устная версия ответа на 5
минут}\label{ux443ux441ux442ux43dux430ux44f-ux432ux435ux440ux441ux438ux44f-ux43eux442ux432ux435ux442ux430-ux43dux430-5-ux43cux438ux43dux443ux442}

Сначала я формулирую главную теорему. Пусть \(X:\Omega\to E\) -
случайный элемент, а \(P_X\) - его распределение:

\[
P_X(B)=P\{X\in B\}.
\]

Тогда для измеримой функции \(g\):

\[
\mathbb E g(X)
=
\int_\Omega g(X(\omega))\,dP
=
\int_E g(x)\,P_X(dx),
\]

если \(g\ge0\) или если \(g(X)\) интегрируема. Это называется заменой
переменной в интеграле Лебега через образ меры.

Дальше говорю: если \(X\) вещественная, можно писать

\[
\mathbb E g(X)=\int_{\mathbb R} g(x)\,dF_X(x).
\]

Если есть плотность \(p_X\), то

\[
\mathbb E g(X)=\int_{\mathbb R} g(x)p_X(x)\,dx.
\]

Важно: плотность существует не всегда. Если плотности нет, надо
интегрировать по \(P_X\), а не по \(p_Xdx\).

Затем формулирую классическую замену переменной. Если \(\Phi:U\to V\) -
\(C^1\)-диффеоморфизм, то

\[
\int_V h(y)\,dy
=
\int_U h(\Phi(x))|\det D\Phi(x)|\,dx.
\]

Отсюда для \(Y=\Phi(X)\):

\[
p_Y(y)=p_X(\Phi^{-1}(y))|\det D\Phi^{-1}(y)|.
\]

После этого перехожу к моментам. Для вещественной случайной величины:

\[
\mathbb EX^k=\int x^k\,P_X(dx),
\qquad
\mathbb E|X|^k=\int |x|^k\,P_X(dx).
\]

При плотности:

\[
\mathbb EX^k=\int x^kp_X(x)\,dx.
\]

Дисперсия:

\[
\operatorname{Var}X
=
\int (x-m)^2\,P_X(dx)
=
\mathbb EX^2-m^2.
\]

Для пары:

\[
\operatorname{Cov}(X,Y)
=
\int (x-m_X)(y-m_Y)\,P_{X,Y}(dx,dy).
\]

Для вектора:

\[
\Sigma
=
\int (x-m)(x-m)^T\,P_X(dx).
\]

И завершаю свойством:

\[
a^T\Sigma a=\operatorname{Var}(a^TX)\ge0.
\]

Главные ловушки: не писать плотность без абсолютной непрерывности, не
забывать существование моментов и якобиан при классической замене
переменной.

\subsection{15. Если осталось 2 минуты перед
экзаменом}\label{ux435ux441ux43bux438-ux43eux441ux442ux430ux43bux43eux441ux44c-2-ux43cux438ux43dux443ux442ux44b-ux43fux435ux440ux435ux434-ux44dux43aux437ux430ux43cux435ux43dux43eux43c}

Запомнить четыре блока.

\textbf{Замена переменной через распределение:}

\[
\mathbb E g(X)
=
\int_\Omega g(X)\,dP
=
\int_E g(x)\,P_X(dx).
\]

\textbf{Если есть плотность:}

\[
\mathbb E g(X)
=
\int g(x)p_X(x)\,dx.
\]

\textbf{Моменты и дисперсия:}

\[
\mathbb EX^k=\int x^k\,P_X(dx),
\qquad
\operatorname{Var}X=\int (x-m)^2\,P_X(dx).
\]

\textbf{Векторный случай:}

\[
m=\int x\,P_X(dx),
\qquad
\Sigma=\int (x-m)(x-m)^T\,P_X(dx).
\]

Главная ловушка: интеграл по плотности допустим только при наличии
плотности; иначе нужно интегрировать по мере распределения \(P_X\).

\newpage

\end{document}
